\documentclass[10pt]{article}

\usepackage[utf8]{inputenc}

\usepackage[T1]{fontenc}

\usepackage{amsmath}

\usepackage{amsfonts}

\usepackage{amssymb}

\usepackage[version=4]{mhchem}

\usepackage{stmaryrd}

\usepackage{bbold}



\begin{document}

имодействие с внешним классич. током $j_{a}(x)$ описывается добавлением к Л. члена



\[

\sum_{a} j_{a}(x) \varphi^{a}(x)

\]



Принципиальное для квантовой теории поля требование перенормируемости налагает новые жесткие ограничения на вид Л.; в болышинстве реальных моделей остающаяся свобода сводится к выбору неболышого числа констант $\begin{array}{ll}\text { (масс и констант взаимодействия). } & \text { В. П. Павлов. }\end{array}$



ЛАКСА ПАРА - см. Обратной задачи рассеяния метод.



ЛАКСА ПРЕДСТАВЛЕНМЕ - представление динамической системы



\[

\dot{u}_{i}=f_{i}\left(u_{1}, \ldots, u_{n}\right), i=1,2, \ldots, n,

\]



для которой сушествуют такие матрицы $L, A$ (элементы которых являются функциями от $u_{1}, \ldots, u_{n}$ ), что система (1) эквиватентна уравнению Лакса



\[

\dot{L}=[L, A]=L A-A L .

\]



В этом случае собственные числа матрицы $L$ являются первыми интегралами системы (1).



Л. п. содержит спектральный параметр $E$, если матрицы $L(E)$ и $A(E)$ разлагаются в ряд по степеням $E$ и уравнение (2) справедливо в силу системы (1) при всех значениях параметра $E$. В этом случае все коэффициенты разложения функции $\operatorname{Tr}\left(L(E)^{k}\right), k=1,2, \ldots$, по степеням $E$ являются первыми интегралами системы (1). Во многих случаях из Л. П. со спектральным параметром следует возможность явного интегрирования системы (1) в тета-функциях римановых поверхностей, заданных уравнением



\[

R(E, w)=\operatorname{det}(L(t, E)-w \cdot 1)=0 .

\]



Первое Л. п. со спектральным параметром было указано Ф. Кеттером (F. Kötter, 1900) для интегрируемых случаев В. А. Стеклова и А. М. Ляпунова движения твердого тела в жидкости. В 1974 С. В. Манаков и Г. Флашка (H. Flaschka) нашли Л. П. для системы взаимодействующих частиц на прямой (экспоненциальная модель кристаллич. решетки це почка Тода).



Л. П. со спектральным параметром построено для всех классич. интегрируемых случаев динамики твердого тела, имеющего неподвижную точку, и для ряда новых интегрируемых задач механики и математич. физики.



См. также Синус-Гордона уравнение.



Лит.: [1] Итоги науки и техники. Современные проблемы математики. Динамические системы, т. 7, М., 1987; [2] Лакс П. Д., «Математика», 1969, т. 13, № 5, с. 128-150; [3] Б о го я в ле нс к и й О. И., Опрокидывающиеся солитоны. Нелинейные интегрируемые уравнения, М., 1990.



О. И. Богоявленский.



ЛАКСА - ФИЛЛИПСА ПОДХОД в те о р и и ра с с е я н и я - подход, применяемый к анализу задач рассеяния для унитарных групп эволюции $U_{t}$ в гильбертовом пространстве $H$ с одномерным непрерывным либо дискретным временем $t$ в случае, когда эти группы обладают ортогональными приходящими и уходящими подпространствами $D_{\mp} \subset H$, $D_{-} \perp D_{+}:$



\[

\begin{aligned}

& U_{t} D_{+} \subset D_{+}, t>0 ; U_{t} D_{+} \rightarrow 0, t \rightarrow+\infty, \\

& U_{t} D_{-} \subset D_{-}, t<0 ; U_{t} D_{-} \rightarrow 0, t \rightarrow-\infty .

\end{aligned}

\]



Если две группы эволюции $U_{t}, U_{t}^{0}$, действующие в различных гильбертовых пространствах $H, H^{0}$, совпадают на общих приходящих и уходящих подпространствах $D_{\mp} \subset H \cap H^{0}$, то, используя в качестве оператора отождествления операторы ортогонального проектирования $P_{ \pm}$на $D_{ \pm}$, легко установить существование волновых операторов:



\[

\begin{aligned}

& W_{-}\left(U^{0}, U\right)=\lim _{t \rightarrow-\infty} U_{-t}^{0} P_{-} U_{t}, \\

& W_{+}\left(U^{0}, U\right)=\lim _{t \rightarrow \infty} U_{-t}^{0} P_{+} U_{t} .

\end{aligned}

\]



316 ЛАKCA Если невозмущенная группа $U_{t}^{0}$ является т р и в и а л в н ы м с це плением изометрич. полугрупп $U_{t}^{0} \mid D_{+}, \quad t>0$, $U_{t}^{0} \mid D_{-}, t<0$, то есть



\[

H^{0}=D_{-} \otimes D_{+}

\]



то оператор рассеяния



\[

S=W_{-}\left(U^{0}, U\right) W_{+}^{*}\left(U_{0}, U\right),

\]



коммутирующий с невозмущенной группой, переводит $D_{+}$ в себя. В спектральном представлении невозмущенной группы он реализуется в виде оператора умножения на матричную сжимающую функцию $S(k): E \rightarrow E-$ матрицу рассеяния. Если использовать приходящее спектральное представление, где $H$ реализуется как пространство квадратично-суммируемых вектор-функций на вещественной оси $L_{2}(\mathbb{R}, E)$, а $D_{-}-$в виде класса Харди всех функций из $L_{2}(\mathbb{R}$, $E)$, допускающих аналитич. продолжение в нижнюю полуплоскость, то матрица рассеяния в этом представлении оказывается аналитической в верхней полуплоскости $\operatorname{Im} k>0$. Она унитарна на вещественной оси в том и только в том случае, когда



\[

\bigvee_{t>0} U_{t} D_{-}=V_{t<0} U_{t} D_{+}=H .

\]



В рамках Л.-Ф.П. легко устанавливается связь между асимптотич. поведением оператора эволюции $U_{t}$ при $t \rightarrow \pm \infty$ в зоне рассеяния и поведением его срезок $Z_{t}$ на трансляционно-инвариантное подпространство $K=H \ominus\left(D_{-} \oplus D_{+}\right):$



\[

Z_{t}=P_{k} U_{t} \mid K, t \geqslant 0 .

\]



Из условия ортогональности $D_{-} \perp D_{+}$сразу следует, что семейство срезок $Z_{t}$ образуег полугруппу сжатий. В случае непрерывного времени ее генератор



\[

B=\lim _{t \rightarrow+0}\left(Z_{t}-I\right)(i t)^{-1}

\]



является диссипативным оператором, чья характеристич. функция совпадает с матрицей рассеяния. В терминах характеристич. функции спектральный анализ оператора $B$ может быть проведен в явном виде. В частности, спектр оператора $B$ совпадает с множеством сингулярностей $S$-матрицы на комплексной плоскости переменного $k$ : полюсы $S^{-1}(k)$ суть собственные числа оператора $B$, а разрезы совпадают с непрерывным спектром (см. Неспектральные особенности). Аналогичные факты имеют место и в случае дискретного времени.



Если нек-рая полоса в верхней полуплоскости $0 \leqslant \operatorname{Im} k \leqslant \gamma$ свободна от спектра генератора $B$, отвечающего рассматриваемой задаче рассеяния, то имеет место экспоненциальное убывание энергии в трансляционно-инвариантном подпространстве $K$ :



\[

\left\|P_{K} U_{t} f\right\|=\left\|Z_{t} f\right\| \leqslant e^{-\gamma t}\|f\|, f \in K, t \geqslant 0 .

\]



Это заведомо выполнено, если $S$-матрица унитарна и полугруппа $Z_{t}$ компактна при нек-ром $t$.



Пользуясь описанным подходом, П. Лакс и Р. Филлипс [1] установили экспоненциальное убывание энергии акустич. волн вблизи компактного звездного препятствия $\Omega$ в $\mathbb{R}^{3}$. B этом случае роль группы $U_{t}$ играет разрешающая группа волнового уравнения $u_{t t}=\Delta u,\left.u\right|_{\partial \Omega}=0$, преобразуюшая данные Коши $u=\left(u_{0}, u_{1}\right)$ от момента времени $t=0$ к моменту времени $t$. Эта группа унитарна в энергетич. норме $\|u\|_{\varepsilon}$ :



\[

\|u\|_{\varepsilon}^{2}=\frac{1}{2} \int_{\mathbb{R}^{3}}\left(\left|\nabla u_{0}\right|^{2}+\left|u_{1}\right|^{2}\right) d x

\]



Роль приходящих и уходящих подпространств играют данные Коши приходящих и уходящих волн, сосредоточенные внутри шара $K_{a}=\{x:|x|<a\}$, содержащего препятствие:



\[

\begin{aligned}

& D_{-}=\left\{\left(u_{0}, u_{1}\right): u(x, t)=0,|x|+t<a\right\}, \\

& D_{+}=\left\{\left(u_{0}, u_{1}\right): u(x, t)=0,|x|-t<a\right\},

\end{aligned}

\]



а роль трансляционно-инвариантного подпространства $K$ - семейство данных Коши с ограниченной энергией:





\end{document}